\subsection{基域扩张}

设 \(\mathfrak{g}\) 是一个 K 上的李代数，\(\mathbf{K}_1\) 是 K 的扩域。显然，\(\mathfrak{g}_{(\mathbf{K}_1)}\) 可解，当且仅当 \(\mathfrak{g}\) 可解，因为 \(\mathcal{D}^n (\mathfrak{g}_{(\mathbf{K}_1)}) = (\mathcal{D}^n\mathfrak{g})_{(\mathbf{K}_1)}\)。

设 \(\mathfrak{r}\) 是 \(\mathfrak{g}\) 的根基。则 \(\mathfrak{r}_{(\mathbf{K}_1)}\) 是 \(\mathfrak{g}_{(\mathbf{K}_1)}\) 的根基。事实上，设 \(\beta\) 是 \(\mathfrak{g}\) 的 Killing 形式。由于 \(\mathfrak{r}\) 是 \(\mathcal{D}\mathfrak{g}\) 关于 \(\beta\) 的正交子空间（命题 5 (b)），\(\mathfrak{r}_{(\mathbf{K}_1)}\) 是 \((\mathcal{D}\mathfrak{g})_{(\mathbf{K}_1)} = \mathcal{D}(\mathfrak{g}_{(\mathbf{K}_1)})\) 关于由 \(\beta\) 从 \(\mathbf{K}\) 扩张到 \(\mathbf{K}_1\) 所得到的形式的正交子空间；该形式就是 \(\mathfrak{g}_{(\mathbf{K}_1)}\) 的 Killing 形式（\(\S 3\)，第 8 号）。我们的断言于是由再次应用命题 5 (b) 得到。

